%% kapa: exam=1 round=19 year=2008 period=1 form=A subject=accounting subject_ko=회계학 seq=15 number=95
%% source: Q-Net 감정평가사 기출문제 (cst00309) artlSeq=1001962 file=extracted/1008/1001962/2008년 제19회 감정평가사 1차시험 기출문제/1교시.hwp
%% conversion: kordoc 4.24.51 -> tools/gen1.py (auto); review=proofread-sonnet; flags=none
\begin{question}{95}
재고자산과 관련된 다음의 설명 중 적절하지 않은 것은?
\begin{choices}
  \item 후입선출법을 사용하여 재고자산의 원가를 결정한 경우에는 대차대조표가액과, 선입선출법 또는 평균법에 저가법을 적용하여 계산한 재고자산평가액과의 차이를 주석으로 기재한다.
  \item 저가기준을 적용하여 매출가격환원법을 사용하는 경우에는 원가율을 계산할 때 가격인하를 매출가격에 의한 판매가능액에서 차감하지 아니한다.
  \item 재료원가, 노무원가 및 기타의 제조원가 중 비정상적으로 낭비된 부분은 재고자산 원가에 포함할 수 없으며 발생기간의 비용으로 인식하여야 한다.
  \item 재고자산의 장부상 수량과 실제 수량과의 차이에서 발생하는 감모손실은 영업외비용으로 분류한다.
  \item 재고자산 평가를 위한 저가법은 종목별로 적용한다. 그러나 재고항목들이 서로 유사하거나 관련되어 있는 경우에는 저가법을 조별로 적용할 수 있다.
\end{choices}
\end{question}
\answer{4}
